\documentclass[UTF8,a4paper,zihao=-4]{ctexart}

\usepackage[top=2.4cm,bottom=2.3cm,left=2.35cm,right=2.35cm]{geometry}
\usepackage{fontspec}
\usepackage{booktabs}
\usepackage{tabularx}
\usepackage{longtable}
\usepackage{array}
\usepackage{enumitem}
\usepackage[table]{xcolor}
\usepackage{hyperref}
\usepackage{fancyhdr}
\usepackage{tikz}
\usetikzlibrary{arrows.meta,positioning,fit,backgrounds}

\setCJKmainfont{Songti SC}
\setCJKsansfont{PingFang SC}
\setCJKmonofont{Kaiti SC}
\setmainfont{Times New Roman}
\setsansfont{Arial}
\setmonofont{Menlo}

\definecolor{Primary}{HTML}{153B5B}
\definecolor{Secondary}{HTML}{167D8D}
\definecolor{Accent}{HTML}{D97834}
\definecolor{LightBlue}{HTML}{EAF3F7}
\definecolor{LightGray}{HTML}{F4F6F8}
\definecolor{RuleGray}{HTML}{CBD5DC}
\definecolor{Success}{HTML}{147D64}

\hypersetup{
  colorlinks=true,
  linkcolor=Primary,
  urlcolor=Secondary,
  citecolor=Primary,
  pdftitle={derp-rs 高性能 Rust DERP 中继服务器项目报告},
  pdfauthor={郑则燊},
  pdfsubject={DERP server design, compatibility and performance evaluation}
}

\setlength{\parindent}{2em}
\setlength{\parskip}{0.22em}
\setlength{\headheight}{24pt}
\setlength{\emergencystretch}{2em}
\setlength{\tabcolsep}{5pt}
\renewcommand{\arraystretch}{1.32}
\raggedbottom
\setlist[itemize]{leftmargin=2em,itemsep=0.25em,topsep=0.3em}
\setlist[enumerate]{leftmargin=2.2em,itemsep=0.25em,topsep=0.3em}

\pagestyle{fancy}
\fancyhf{}
\fancyhead[L]{\small\color{Primary}derp-rs 项目报告}
\fancyhead[R]{\small\color{Primary}郑则燊 \quad 24344151}
\fancyfoot[C]{\small\color{gray}\thepage}
\renewcommand{\headrulewidth}{0.4pt}
\renewcommand{\headrule}{\hbox to\headwidth{\color{RuleGray}\leaders\hrule height \headrulewidth\hfill}}

\ctexset{
  section = {
    format = \Large\bfseries\sffamily\color{Primary},
    beforeskip = 1.0em,
    afterskip = 0.55em
  },
  subsection = {
    format = \large\bfseries\sffamily\color{Secondary},
    beforeskip = 0.85em,
    afterskip = 0.4em
  },
  subsubsection = {
    format = \normalsize\bfseries\sffamily\color{Primary}
  }
}

\newcolumntype{Y}{>{\raggedright\arraybackslash}X}
\newcolumntype{C}[1]{>{\centering\arraybackslash}p{#1}}
\newcolumntype{R}[1]{>{\raggedleft\arraybackslash}p{#1}}
\newcommand{\code}[1]{\texttt{\small #1}}
\newcommand{\status}{\textcolor{Success}{\textbf{完整}}}
\newcommand{\metric}[1]{\textcolor{Primary}{\textbf{#1}}}

\begin{document}

\begin{titlepage}
  \thispagestyle{empty}
  \vspace*{0.7cm}
  \noindent{\color{Primary}\rule{\textwidth}{2.2pt}}

  \vspace{1.1cm}
  {\sffamily\bfseries\fontsize{30}{39}\selectfont\color{Primary}
    derp-rs\par}
  \vspace{0.35cm}
  {\sffamily\bfseries\fontsize{21}{30}\selectfont
    高性能 Rust DERP 中继服务器\par}
  \vspace{0.7cm}
  {\Large\color{Secondary}项目设计、协议兼容与性能评估报告\par}

  \vspace{1.3cm}

  \begin{center}
    \renewcommand{\arraystretch}{1.6}
    \begin{tabular}{>{\sffamily\bfseries\color{Primary}}m{3.3cm}m{8.7cm}}
      项目名称 & derp-rs 高性能 Rust DERP 中继服务器 \\
      项目成员 & 郑则燊 \\
      学号 & 24344151 \\
      项目类型 & 网络系统软件 / 开源协议兼容实现 \\
      主要语言 & Rust 2024 Edition \\
      对比基线 & Tailscale derper v1.100.0 \\
      报告日期 & 2026 年 7 月 \\
      项目仓库 & \url{github.com/jaisonZheng/derp-rs} \\
    \end{tabular}
  \end{center}
\end{titlepage}

\pagenumbering{Roman}

\section*{摘要}
\addcontentsline{toc}{section}{摘要}

DERP（Designated Encrypted Relay for Packets）是 Tailscale 在节点无法建立直接
WireGuard 通道时使用的中继协议。本项目以公开协议、官方源码和测试为兼容基线，
使用 Rust 从零实现生产可用的高性能 DERP 服务器。系统覆盖 DERP v2 加密握手、
全部稳定帧类型、Fast Start、标准 HTTP Upgrade、WebSocket、STUN、区域 Mesh、
重复连接、反向路径通知、准入控制、限速、背压、TLS、可观测性和优雅重启。

实现采用 Tokio 异步运行时、DashMap 活跃连接索引、有界 MPSC 队列、
\code{bytes::Bytes} 引用计数载荷和约 64 KiB 的批量写缓冲。兼容性方面，将
Tailscale 官方 Go 测试中可观察的线缆行为改写为外部服务器黑盒测试，并使用
官方 Go 客户端 v1.100.0 验证；本机、GitHub Actions 和腾讯云 Linux 环境均为
12/12 通过。

性能实验显示，在 Apple M3 loopback、16 个官方客户端、1,200 字节载荷的
五轮测试中，derp-rs 中位吞吐为 366,442 packets/s，官方 Go derper 为
263,363 packets/s，前者高 39.14\%。在腾讯云 2 vCPU Linux 主机上，针对
100 到 10,000 条连接以及明文、TLS、持续转发、慢接收端和连接抖动等工作负载，
derp-rs 在所有已测场景中的稳态与峰值 RSS 均显著低于官方实现。

\textbf{关键词：} DERP；Rust；Tailscale；异步网络；协议兼容；高并发；内存优化

\clearpage
\begingroup
\footnotesize
\tableofcontents
\endgroup
\clearpage
\pagenumbering{arabic}

\section{项目背景与目标}

\subsection{项目背景}

WireGuard 节点通常优先通过 UDP 建立端到端直连，但对称 NAT、严格防火墙、
运营商网络和受限企业网络可能使打洞失败。DERP 在这种情况下提供可靠的 TCP/TLS
中继路径。中继服务器只依据 NodeKey 转发已经端到端加密的数据包，不能解密
WireGuard 业务载荷，因此它既是可用性基础设施，也是高连接密度的网络服务。

官方 derper 由 Go 实现，功能成熟、生态完整。本项目选择 Rust 的原因不是简单
重写，而是验证以下工程命题：在保持协议兼容和生产功能的前提下，是否可以利用
Rust 的所有权模型、紧凑数据结构、无垃圾回收运行时和显式背压，获得更高吞吐和
更稳定的内存占用。

\subsection{项目目标}

\begin{enumerate}
  \item \textbf{协议完整性：}实现 DERP v2 加密握手与全部稳定帧
    \code{0x01} 到 \code{0x15}，兼容官方 Tailscale 客户端。
  \item \textbf{生产功能：}支持 TLS、WebSocket、STUN、区域 Mesh、准入控制、
    速率限制、有界队列、Prometheus 指标与优雅重启。
  \item \textbf{高性能：}降低热路径分配、复制、锁竞争和系统调用次数，提高
    relay throughput。
  \item \textbf{低内存：}在 100 到 10,000 条连接、慢客户端和 churn 场景下，
    保持可预测的稳态与峰值 RSS。
  \item \textbf{可验证：}使用官方 Go 客户端和公开测试行为进行黑盒互操作测试，
    保留可重复运行的基准脚本与机器可读原始数据。
\end{enumerate}

\subsection{范围与边界}

本项目覆盖 DERP 线协议和服务器核心行为。官方 \code{cmd/derper} 中的 ACME/GCP
证书管理、本机 \code{tailscaled} 调用、Linux XDP 快速路径和实验性 ACE proxy
属于平台集成或实验功能，不是稳定 DERP v2 互操作的必要条件。derp-rs 使用 PEM
证书、通用 HTTP Admission Controller 和静态 bootstrap DNS 文件替代相应平台绑定。

\section{官方功能覆盖}

\subsection{核心能力矩阵}

\rowcolors{2}{LightGray}{white}
\begin{longtable}{
  >{\raggedright\arraybackslash}p{0.18\textwidth}
  >{\raggedright\arraybackslash}p{0.60\textwidth}
  C{0.12\textwidth}}
\toprule
\rowcolor{Primary}
\color{white}\textbf{模块} &
\color{white}\textbf{已实现的官方行为} &
\color{white}\textbf{状态} \\
\midrule
\endfirsthead
\toprule
\rowcolor{Primary}
\color{white}\textbf{模块} &
\color{white}\textbf{已实现的官方行为} &
\color{white}\textbf{状态} \\
\midrule
\endhead
连接入口 &
原生 \code{Upgrade: DERP}、Fast Start 与 WebSocket \code{derp} 子协议 &
\status \\
加密身份 &
ServerKey、NodeKey、NaCl \code{crypto\_box}、ClientInfo 与 ServerInfo &
\status \\
DERP v2 转发 &
NodeKey 寻址、接收包源 NodeKey、64 KiB 最大载荷、未知帧安全跳过 &
\status \\
反向路径 &
源节点最后一条连接断开时向既有接收者发送 PeerGone(Disconnected) &
\status \\
目的端缺失 &
普通包静默，合法 disco wrapper 返回 PeerGone(NotHere)，每连接 3 次/秒 &
\status \\
重复连接 &
同 NodeKey 多会话、Health 状态、最新活动会话选择与恢复通知 &
\status \\
控制帧 &
Ping、Pong、KeepAlive、NotePreferred、Health 与 Restarting &
\status \\
区域 Mesh &
PSK、WatchConns、PeerPresent/Gone、ForwardPacket 与 ClosePeer &
\status \\
STUN &
Tailscale Binding Request 校验与 IPv4/IPv6 XOR-MAPPED-ADDRESS &
\status \\
HTTP 运维 &
probe、latency-check、generate\_204、bootstrap-dns、metrics 与 debug &
\status \\
安全和资源 &
PEM TLS、Admission Controller、令牌桶、有界队列和写超时 &
\status \\
\bottomrule
\end{longtable}
\rowcolors{2}{}{}

\clearpage
\subsection{帧类型覆盖}

\centering
\small
\begin{tabular}{clclcl}
\toprule
\textbf{值} & \textbf{类型} & \textbf{值} & \textbf{类型} &
\textbf{值} & \textbf{类型} \\
\midrule
\code{0x01} & ServerKey      & \code{0x06} & KeepAlive     & \code{0x11} & ClosePeer \\
\code{0x02} & ClientInfo     & \code{0x07} & NotePreferred & \code{0x12} & Ping \\
\code{0x03} & ServerInfo     & \code{0x08} & PeerGone      & \code{0x13} & Pong \\
\code{0x04} & SendPacket     & \code{0x09} & PeerPresent   & \code{0x14} & Health \\
\code{0x05} & RecvPacket     & \code{0x0a} & ForwardPacket & \code{0x15} & Restarting \\
              &                & \code{0x10} & WatchConns     &              & \\
\bottomrule
\end{tabular}
\par\smallskip
表 2：DERP 稳定帧覆盖
\par
\normalsize
\raggedright

\section{系统设计与实现思路}

\subsection{总体架构}

\begin{figure}[htbp]
\centering
\begin{tikzpicture}[
  box/.style={
    draw=Primary, rounded corners=2pt, fill=LightBlue,
    minimum height=10mm, align=center, font=\small\sffamily
  },
  core/.style={
    draw=Secondary, rounded corners=2pt, fill=Secondary!10,
    minimum height=11mm, align=center, font=\small\sffamily\bfseries
  },
  ext/.style={
    draw=Accent, rounded corners=2pt, fill=Accent!8,
    minimum height=10mm, align=center, font=\small\sffamily
  },
  arrow/.style={-{Latex[length=2mm]}, thick, draw=Primary}
]
  \node[box, text width=36mm] (clients) at (2.2,5.0)
    {Tailscale / Headscale\\客户端};
  \node[box, text width=36mm] (mesh) at (7.3,5.0)
    {区域 Mesh\\官方或 Rust derper};
  \node[ext, text width=28mm] (stunclient) at (12.0,5.0)
    {STUN 客户端};

  \node[box, text width=40mm] (transport) at (2.2,3.25)
    {TCP / TLS / WebSocket\\HTTP Upgrade 与 Fast Start};
  \node[box, text width=40mm] (handshake) at (7.3,3.25)
    {加密握手与准入\\NodeKey + crypto\_box};
  \node[ext, text width=30mm] (stun) at (12.0,3.25)
    {UDP STUN\\XOR-MAPPED-ADDRESS};

  \node[core, text width=40mm] (relay) at (2.2,0.70)
    {Relay Core\\活跃连接与反向路径索引};
  \node[core, text width=40mm] (session) at (7.3,0.70)
    {Session Tasks\\有界队列与批量写缓冲};
  \node[ext, text width=30mm] (ops) at (12.0,0.70)
    {Metrics / Debug\\Admission HTTP};

  \draw[arrow] (clients) -- (transport);
  \draw[arrow] (mesh) -- (handshake);
  \draw[arrow] (stunclient) -- (stun);
  \draw[arrow] (transport) -- (relay);
  \draw[arrow] (handshake) -- (session);
  \draw[arrow, <->] (relay) -- (session);
  \draw[arrow] (stun) -- (ops);
  \draw[arrow, <->] (session) -- (ops);
\end{tikzpicture}
\caption{derp-rs 总体架构}
\end{figure}

TCP 主监听器接受连接后，根据配置可先执行 rustls TLS 握手。HTTP 层识别
Fast Start、标准 DERP Upgrade 或 WebSocket，并将统一的异步字节流交给 DERP
会话层。UDP STUN 使用独立任务，不占用 DERP TCP 会话资源。Metrics、debug
和 admission 复用轻量 HTTP 处理逻辑。

\subsection{连接与加密握手}

\begin{enumerate}
  \item 服务器发送包含固定 magic 和 32 字节公钥的 ServerKey 帧。
  \item 客户端生成或加载 NodeKey，将 ClientInfo 使用双方 Curve25519 密钥和
    24 字节 nonce 加密后发送。
  \item 服务器解密并校验协议版本、Mesh PSK 和客户端能力；如配置 admission，
    则把 NodePublic 与来源地址提交给控制器。
  \item 服务器注册会话，返回加密 ServerInfo，之后进入稳态 frame 循环。
\end{enumerate}

NodeKey 是路由地址而不是业务内容密钥。DERP 只处理外层已加密数据包，无法读取
内部 WireGuard 流量。Mesh PSK 使用常量时间比较，持久服务器私钥首次创建时采用
原子写入和 \code{0600} 权限。

\subsection{高性能转发热路径}

\begin{itemize}
  \item \textbf{双层索引：}DashMap 保存 NodeKey 到当前活跃 Session 的热路径映射；
    同 key 多会话的控制面使用小型锁保护，避免每个数据包都遍历连接集合。
  \item \textbf{少复制载荷：}数据包进入服务端后使用 \code{bytes::Bytes}，
    本地转发与 Mesh 尝试共享引用计数内存。
  \item \textbf{显式背压：}每个客户端使用固定深度的 MPSC 队列；满队列立即计数并
    丢弃，慢客户端不会让服务器内存无限增长。
  \item \textbf{批量写：}写任务将多个 frame 直接编码到同一个约 64 KiB 缓冲区，
    减少临时 payload、二次复制和小系统调用。
  \item \textbf{锁与计数：}Prometheus 指标采用原子计数；本地转发查找不持有
    全局互斥锁。
\end{itemize}

\subsection{连接生命周期与内存优化}

原始 Rust 版本已经比 Go 基线节省内存，但高连接数实验显示仍可进一步优化。
最终实现做了以下调整：

\begin{itemize}
  \item 无数据的空闲连接不分配批量写缓冲；
  \item 多帧直接编码到单一缓冲，不再先生成临时 payload 再复制；
  \item 写缓冲按观测批次增长，连续 15 秒无数据后释放过大容量；
  \item 单批约 64 KiB，避免聚合导致不可控延迟和峰值；
  \item 删除两个冗余的 per-session 原子堆分配；
  \item 常见的单会话 PeerSet 使用 inline small-vector 存储。
\end{itemize}

\subsection{Mesh 与故障传播}

受信任的区域节点通过加密 ClientInfo 携带 Mesh PSK，之后使用 WatchConns 订阅
PeerPresent 与 PeerGone。目标节点不在本机时，Relay Core 查询 Mesh 路由并发送
ForwardPacket。已经转发过的包不会再次进入 Mesh，避免环路。成功的 A 到 B 转发
记录反向路径；A 最后一条连接离开后，B 收到 PeerGone(Disconnected)。

\section{测试设计}

\subsection{多层验证策略}

\begin{table}[htbp]
\centering
\small
\begin{tabularx}{\textwidth}{p{0.20\textwidth}Y Y}
\toprule
\textbf{层级} & \textbf{主要内容} & \textbf{作用} \\
\midrule
Rust 单元测试 &
frame round-trip、官方 greeting 向量、crypto box、路由、重复连接、限速、STUN &
快速验证纯逻辑和边界条件 \\
Rust 集成测试 &
WebSocket 握手与 Ping/Pong、Rust 到 Rust regional mesh &
验证异步服务组合 \\
官方客户端黑盒 &
官方 Go DERP HTTP 客户端连接真实 release 进程 &
避免“服务端和自制客户端同时写错” \\
社区行为适配 &
双向转发、最新连接路由、非 Fast Start Upgrade &
补充独立 Rust 社区实现视角 \\
生产压力测试 &
100 到 10,000 连接、TLS、active、slow、churn &
验证资源边界和内存峰值 \\
\bottomrule
\end{tabularx}
\caption{验证层级}
\end{table}

Tailscale 上游测试通常把 Go \code{derpserver.Server} 直接嵌入测试进程，不能
直接指向第三方二进制。本项目将可观察的协议行为改写为外部服务器测试，官方
Go DERP HTTP 客户端固定为 v1.100.0。Go 内部数据结构、XDP 实现细节和只测试
客户端错误处理的用例不做机械移植。

\subsection{黑盒一致性结果}

\begin{table}[htbp]
\centering
\begin{tabularx}{0.90\textwidth}{Y C{2.7cm}}
\toprule
\textbf{环境} & \textbf{结果} \\
\midrule
Apple Silicon macOS，本机 release 进程 & \metric{12/12 通过} \\
腾讯云 Ubuntu 24.04 x86\_64，CI release ELF & \metric{12/12 通过} \\
GitHub Actions Ubuntu，Rust 测试 + Clippy + 黑盒套件 & \metric{全部通过} \\
\bottomrule
\end{tabularx}
\caption{跨平台一致性验证}
\end{table}

黑盒套件覆盖三节点转发、双向转发、源 NodeKey、Ping/Pong、反向 PeerGone、
NotHere disco 规则和 burst 限速、Mesh watcher 快照与上下线、重复连接 Health、
最新连接路由、NotePreferred、Fast Start、标准 HTTP 101 Upgrade、运维端点和
官方 STUN 请求/解析。

\section{实验结果}

\subsection{吞吐实验方法}

\begin{itemize}
  \item 硬件：Apple M3，macOS 15.6，arm64。
  \item 服务端：derp-rs 0.1.0 release 与官方 derper v1.100.0。
  \item 客户端：16 个官方 Go DERP 客户端组成转发环。
  \item 负载：1,200 字节 payload，每轮 256,000 个确认送达的数据包。
  \item 传输：loopback TCP、Fast Start、关闭 TLS 与 STUN。
  \item 统计：每个实现运行 5 次，报告 packets/s 中位数。
\end{itemize}

负载生成器只有在接收端验证所有发送包后才推进窗口，因此结果不能通过静默丢包
获得。该实验强调 framing、路由、分配、队列、调度和 socket write，不代表公网 RTT。

\subsection{吞吐结果}

\begin{table}[htbp]
\centering
\small
\begin{tabularx}{\textwidth}{p{0.25\textwidth}Y R{2.3cm}R{2.2cm}}
\toprule
\textbf{服务端} & \textbf{五轮 packets/s} &
\textbf{中位数} & \textbf{有效载荷} \\
\midrule
derp-rs 0.1.0 &
380,213；366,442；375,886；289,351；312,408 &
\metric{366,442} & \metric{3.518 Gbit/s} \\
Go derper v1.100.0 &
259,203；280,457；256,018；263,363；286,846 &
263,363 & 2.528 Gbit/s \\
\bottomrule
\end{tabularx}
\caption{同机五轮 DERP 转发吞吐}
\end{table}

derp-rs 的中位吞吐比官方实现高 \metric{39.14\%}。同一微基准中的 RSS 样本为
4,704 KiB 对 22,208 KiB，Rust 低 78.82\%，Go 进程为 Rust 的 4.72 倍。
RSS 样本不是峰值，因此更强的内存结论来自下一节独立的 Linux 生产规模矩阵。

\subsection{生产规模 RSS 实验方法}

实验运行于腾讯云 Ubuntu、Linux 6.8 x86\_64、2 vCPU、1.92 GiB RAM 主机。
服务器每个场景重新启动，客户端使用官方 \code{tailscale.com/derp v1.100.0}：

\begin{itemize}
  \item \textbf{idle：}完成握手并持续读取控制帧；
  \item \textbf{active：}客户端组成环，以总计约 10,000 packets/s 转发并验证接收；
  \item \textbf{slow：}一半接收端停止读取且 socket receive buffer 为 4 KiB；
  \item \textbf{churn：}每 500 ms 关闭并替换约 10\% 活跃连接。
\end{itemize}

预热后，baseline 为 10 次 100 ms 采样均值，steady 为状态就绪后的 20 次
100 ms 采样均值，peak 每 50 ms 采样一次。Linux 数据源为
\code{/proc/<pid>/status} 的 \code{VmRSS}。

\subsection{生产规模 RSS 结果}

\begin{table}[htbp]
\centering
\small
\begin{tabular}{llrrrr}
\toprule
\textbf{传输} & \textbf{工作负载} & \textbf{连接} &
\textbf{Rust/MiB} & \textbf{Go/MiB} & \textbf{降幅} \\
\midrule
明文 & Idle          & 1,000  & \metric{15.9}  & 51.5  & 69.17\% \\
明文 & Idle          & 10,000 & \metric{111.4} & 358.9 & 68.96\% \\
明文 & Active        & 5,000  & \metric{70.0}  & 204.1 & 65.72\% \\
明文 & Slow receiver & 5,000  & \metric{111.9} & 308.1 & 63.69\% \\
明文 & Churn         & 5,000  & \metric{58.3}  & 202.9 & 71.26\% \\
\midrule
TLS & Idle           & 1,000  & \metric{24.0}  & 79.1  & 69.67\% \\
TLS & Idle           & 10,000 & \metric{184.9} & 595.1 & 68.93\% \\
TLS & Active         & 5,000  & \metric{107.3} & 324.7 & 66.95\% \\
TLS & Slow receiver  & 5,000  & \metric{156.8} & 439.2 & 64.30\% \\
TLS & Churn          & 5,000  & \metric{96.0}  & 314.4 & 69.47\% \\
\bottomrule
\end{tabular}
\caption{选定场景的稳态 RSS}
\end{table}

完整明文矩阵中，derp-rs 稳态 RSS 低 \metric{62.38\% 到 71.36\%}，峰值低
\metric{62.38\% 到 79.85\%}；完整 TLS 矩阵中，稳态低
\metric{63.65\% 到 70.99\%}，峰值低 \metric{63.65\% 到 78.91\%}。

5,000 连接 churn 能体现垃圾回收和连接生命周期造成的峰值差异。明文 peak RSS
为 Rust 58.6 MiB、Go 290.8 MiB；TLS 为 Rust 100.1 MiB、Go 391.4 MiB。
慢接收端场景中有界队列有效限制了无法及时写出的 payload 数量。

\subsection{优化前后效果}

\begin{table}[htbp]
\centering
\small
\begin{tabular}{llrrr}
\toprule
\textbf{传输} & \textbf{5,000 连接} &
\textbf{原始 Rust/MiB} & \textbf{优化 Rust/MiB} & \textbf{降幅} \\
\midrule
明文 & Idle          & 78.0  & \metric{58.6}  & 24.84\% \\
明文 & Active        & 90.2  & \metric{70.0}  & 22.40\% \\
明文 & Slow receiver & 134.4 & \metric{111.9} & 16.74\% \\
TLS & Idle           & 115.6 & \metric{95.6}  & 17.30\% \\
TLS & Active         & 127.0 & \metric{107.3} & 15.51\% \\
TLS & Slow receiver  & 172.3 & \metric{156.8} & 9.01\% \\
\bottomrule
\end{tabular}
\caption{内存优化前后稳态 RSS}
\end{table}

同一批优化还将本地五轮吞吐中位数从 313,906 提升到 366,442 packets/s，
说明内存下降不是通过降低转发工作量换取的。

\subsection{实验结论边界}

实验支持以下结论：在指定 Linux 主机和 100 到 10,000 条连接的测试包络内，
包括 TLS、持续流量、背压和 churn，derp-rs 的稳态与峰值 RSS 始终显著低于
官方 derper v1.100.0。

该结论不是对所有操作系统、分配器、CPU 架构、未来 Go/Rust 版本、扩展配置、
任意流量分布或无限连接数的数学保证。公网性能还会受到带宽、丢包、拥塞和 RTT
主导。生产容量规划应在目标主机重跑仓库脚本。

\section{开源项目参考说明}

\subsection{Tailscale}

本项目最主要的兼容参考是 Tailscale 公开仓库中的 DERP 协议、Go 客户端、
\code{derpserver}、\code{cmd/derper} 与 STUN 实现。开发时检查的主线提交如下，
稳定互操作和性能基线固定为 v1.100.0。
\begin{center}
\code{cfd101f9d773695def27a5f6289fc25ac36ac991}
\end{center}
测试参考包括：

\begin{itemize}
  \item \code{derp/derp\_test.go}：SendRecv、Watch、Ping/Pong、PeerGone；
  \item \code{derp/derphttp/derphttp\_test.go}：官方 HTTP 客户端与 probe；
  \item \code{derp/derpserver/derpserver\_test.go}：重复连接、Mesh、限速；
  \item \code{net/stun/stun\_test.go}：STUN request/response 向量。
\end{itemize}

derp-rs 是依据公开协议和可观察行为编写的独立实现，不包含 Tailscale 商标授权，
也不自称 Tailscale 官方产品。项目采用与上游兼容的 BSD-3-Clause 许可证。

\subsection{Rust 社区实现}

\code{rajsinghtech/rustscale} 的测试型 DERP server 提供了双向转发、
latest-writer 路由和非 Fast Start Upgrade 测试思路。本项目移植其与官方协议
一致的行为。该社区实现要求新连接关闭同 key 旧连接，而当前 Tailscale 官方行为
是保留重复连接并发送 Health，因此冲突部分以官方语义为准，没有机械照搬。

\subsection{主要 Rust 依赖}

\begin{table}[htbp]
\centering
\small
\begin{tabularx}{\textwidth}{p{0.22\textwidth}Y}
\toprule
\textbf{开源组件} & \textbf{在项目中的用途} \\
\midrule
Tokio & 多线程异步运行时、TCP/UDP、任务、channel、signal 和 timer \\
rustls / tokio-rustls & 无 OpenSSL 运行时依赖的 TLS 服务器 \\
crypto\_box & Curve25519 XSalsa20Poly1305 兼容加密盒 \\
DashMap / parking\_lot & 并发连接索引与轻量控制面锁 \\
bytes & 引用计数的网络载荷及缓冲管理 \\
tokio-tungstenite & WebSocket DERP 子协议 \\
serde / serde\_json & ClientInfo、ServerInfo、admission 和配置数据 \\
clap / tracing & 命令行参数与结构化运行日志 \\
\bottomrule
\end{tabularx}
\caption{主要开源依赖及用途}
\end{table}

\section{总结与展望}

本项目完成了一个协议兼容、功能覆盖完整、可独立部署的 Rust DERP 服务器。
与官方 Go derper 的对比表明，显式内存生命周期、有界背压和批量写设计能够在
不牺牲协议正确性的情况下，同时提升吞吐并显著降低连接密度下的 RSS。

后续工作包括：在更多 CPU 架构和云厂商复现实验；增加长时间故障注入、网络丢包
和模糊测试；对 admission 与 Mesh 管理面开展独立安全审计；根据上游协议演进持续
运行黑盒兼容套件。对生产部署而言，仍应结合目标网络的 RTT、带宽和连接分布进行
容量规划。

\clearpage
\appendix
\section{复现实验命令}

\begin{verbatim}
cargo test --all-targets --locked
cargo clippy --all-targets --locked -- -D warnings
./scripts/official-conformance.sh
./scripts/benchmark.sh
./scripts/rss-benchmark.sh
\end{verbatim}

完整生产矩阵可使用：

\begin{verbatim}
TLS_MODE=1 \
SCENARIOS="idle:100 idle:1000 idle:5000 idle:10000 \
active:1000 active:5000 slow:1000 slow:5000 \
churn:1000 churn:5000" \
./scripts/rss-benchmark.sh
\end{verbatim}

\section{参考资料}

\begin{thebibliography}{9}
\bibitem{tailscale-derp}
Tailscale Inc.,
\textit{DERP protocol package},
\url{https://pkg.go.dev/tailscale.com/derp}.

\bibitem{tailscale-derpserver}
Tailscale Inc.,
\textit{DERP server implementation},
\url{https://github.com/tailscale/tailscale/tree/main/derp/derpserver}.

\bibitem{tailscale-derper}
Tailscale Inc.,
\textit{cmd/derper},
\url{https://github.com/tailscale/tailscale/tree/main/cmd/derper}.

\bibitem{tailscale-stun}
Tailscale Inc.,
\textit{STUN implementation},
\url{https://github.com/tailscale/tailscale/tree/main/net/stun}.

\bibitem{rustscale}
Raj Singh,
\textit{rustscale DERP test server},
\url{https://github.com/rajsinghtech/rustscale/blob/master/crates/derp/src/server.rs}.

\bibitem{wireguard}
Jason A. Donenfeld,
\textit{WireGuard: Next Generation Kernel Network Tunnel},
\url{https://www.wireguard.com/papers/wireguard.pdf}.
\end{thebibliography}

\end{document}
